\section{Zpětnovazební učení}\label{sec:ai-reinforcement}

Při zpětnovazebním učení agent nedostává pro každý stav předem předepsanou správnou akci. Musí jednat, sledovat následky a~postupně zjišťovat, které volby vedou k~dobrému dlouhodobému výsledku.

Mějme množinu stavů \(\mathcal{S}\) a~pro každý stav \(s\in\mathcal{S}\) množinu přípustných akcí \(\mathcal{A}(s)\). V~čase \(t\) se agent nachází ve stavu \(s_t\), zvolí akci \(a_t\in\mathcal{A}(s_t)\) a~prostředí mu vrátí odměnu \(r_{t+1}\in\R\) a~nový stav \(s_{t+1}\). Jednu zkušenost tedy popisuje čtveřice
\[
    (s_t,a_t,r_{t+1},s_{t+1}).
\]

\begin{figure}[ht]
    \centering
    \begin{tikzpicture}[every node/.style={font=\normalsize}]
    \node[task,fill=blue!10,minimum width=34mm,minimum height=14mm]
        (agent) at (0,0) {agent};
    \node[task,fill=orange!15,minimum width=34mm,minimum height=14mm]
        (env) at (7,0) {prostředí};
    \draw[diredge,bend left=18] (agent.north east) to
        node[above]{akce \(a_t\)} (env.north west);
    \draw[diredge,bend left=18] (env.south west) to
        node[below,align=center]{odměna \(r_{t+1}\)\\stav \(s_{t+1}\)}
        (agent.south east);
    \node[above=3mm of agent] {stav \(s_t\)};
\end{tikzpicture}

    \caption{Cyklus interakce agenta s~prostředím.}
    \label{fig:ai-zpetna-vazba}
\end{figure}

Agent obvykle nechce získat pouze nejvyšší okamžitou odměnu. Zajímá ho diskontovaný součet budoucích odměn
\[
    G(t)
    =
    \sum_{i=0}^{\infty}
    \gamma^i r_{t+i+1},
    \qquad
    0\leqslant\gamma<1.
\]
Parametr \(\gamma\) určuje význam budoucnosti. Pro \(\gamma\) blízké nule převažuje okamžitý zisk; pro \(\gamma\) blízké jedné agent více zohledňuje vzdálenější následky.

\emph{Politika} \(\pi\) přiřazuje stavům akce, tedy \(\pi(s)=a\). Cílem je naučit se politiku, která vede k~co nejvyššímu dlouhodobému zisku. Agent přitom řeší dilema mezi využíváním dosud nejlepší známé akce a~průzkumem jiných možností, které se mohou ukázat jako výhodnější.

\subsection{Q-learning}

Algoritmus Q-learning uchovává pro každou přípustnou dvojici \((s,a)\) hodnotu \(Q(s,a)\), která vyjadřuje, jak výhodné se zatím jeví provést akci \(a\) ve stavu \(s\). Po zkušenosti
\((s_t,a_t,r_{t+1},s_{t+1})\) provede aktualizaci
\[
    Q(s_t,a_t)
    \gets
    Q(s_t,a_t)
    +
    \alpha
    \left[
        r_{t+1}
        +
        \gamma\max_a Q(s_{t+1},a)
        -
        Q(s_t,a_t)
    \right],
\]
kde \(\alpha\in(0,1]\) je rychlost učení. Výraz v~hranaté závorce porovnává dosavadní odhad s~novou informací: okamžitou odměnou a~odhadovanou hodnotou nejlepšího pokračování.

Po učení získáme politiku
\[
    \pi(s)=\argmax_{a\in\mathcal{A}(s)}Q(s,a).
\]

\begin{figure}[ht]
    \centering
    \begin{tikzpicture}[x=1.05cm,y=1.05cm,every node/.style={font=\normalsize}]
    \draw[step=1,black!50] (0,0) grid (5,4);
    \fill[black!25] (2,1) rectangle (3,3);
    \node[align=center] at (2.5,2) {zeď};
    \fill[blue!65] (0.5,0.5) circle (3pt);
    \node[below] at (0.5,0.28) {agent};
    \fill[green!55!black] (4,3) rectangle (5,4);
    \node[text=white] at (4.5,3.5) {cíl};
    \fill[red!65] (4,0) rectangle (5,1);
    \node[text=white] at (4.5,0.5) {past};
    \draw[diredge] (0.7,0.5) -- (1.5,0.5);
    \draw[diredge] (0.5,0.7) -- (0.5,1.5);
\end{tikzpicture}

    \caption{Jednoduché prostředí pro zpětnovazební učení.}
    \label{fig:ai-gridworld}
\end{figure}

Program~\ref{lst:ai-qlearning} implementuje Q-learning v~malém mřížkovém světě. Agent dostává za běžný krok odměnu \(-1\), za dosažení cíle \(10\) a~za vstup do pasti \(-10\). S~pravděpodobností \(\varepsilon\) zvolí náhodnou akci, a~tím prostředí prozkoumává.

\lstinputlisting[
    caption={Q-learning v~jednoduchém mřížkovém světě.},
    label={lst:ai-qlearning}
]{07-ai/code/q_learning.py}

\notebox{Q-learning nepotřebuje předem znát pravidla přechodu prostředí. Hodnoty \(Q\) upravuje pouze podle skutečně získaných zkušeností. Pro velké nebo spojité množiny stavů však tabulka nestačí a~hodnotu \(Q(s,a)\) je nutné aproximovat modelem.}
